\section{Intended uses}
\label{sec:applications}

The construction is general.  Any finite multi-party protocol with
public coins, sealed commitments, public reveals, and deterministic
public payoffs can be expressed.  The design choices were guided by a
particular family of intended uses.  This section sketches them.
None of the patterns described here require new machinery on top of
what \Cref{sec:source-language} through \Cref{sec:refinement} provide.
Backend-specific work instantiates the runtime/refinement layer rather
than changing the source language or event graph.

\subsection{Sealed-bid auctions and mechanism design}

A sealed-bid auction~\citep{Vickrey1961,Myerson1981OptimalAuction}
is the canonical commit/reveal protocol: bidders simultaneously
sealed-bid in a commit phase, the bids are then revealed, and the
allocation and payments are public functions of the revealed
bids.  The surrounding mechanism-design and algorithmic-game-theory
background is standard~\citep{MasColellWhinstonGreen1995,
NisanRoughgardenTardosVazirani2007}.  In the framework:
\begin{itemize}\itemsep1pt
\item Each bidder's bid is a $\Commit$ with a guard restricting
  the value to the bid domain (e.g.\ a bounded integer range).
\item The bids are revealed in a single phase; the
  commit/reveal information barrier (\Cref{thm:barrier}) prevents
  any bidder from observing another bidder's bid before placing
  their own.
\item Allocation rules are public payoff expressions over the
  revealed bids; the payoff-vector kernel games then expose these
  as strategic-form games whose equilibrium predicates are the
  standard incentive-compatibility, individual-rationality, and
  revenue notions.
\item Backend refinement (\Cref{sec:refinement}) is the formal
  link to a contract or auction-house implementation.  A
  stochastic-refinement proof, relative to the backend's chosen
  observation model, rules out implementation traces in which a sealed
  bid becomes observable before the reveal phase.  This is the kind of
  obligation needed to address observation-before-reveal attacks in a
  surrounding execution environment~\citep{Daian2020Flash}.
\end{itemize}

\subsection{Contract runtime backends}

A contract backend adds a machine whose state contains runtime storage
and whose events model the relevant transaction, call, queue, or
block-level actions.  It proves:
\begin{itemize}\itemsep1pt
\item a projection from contract states and observations to the compiled
  event-graph machine;
\item preservation of event and event-batch probability laws under that
  projection;
\item terminal outcome and utility preservation;
\item a strategy-indexed trace law connecting the backend execution
  model to the chosen frontier game surface.
\end{itemize}
The runtime trace-adequacy theorem (\Cref{thm:runtime-trace}) is what
this checklist instantiates.  Those obligations are exactly what the
runtime refinement API isolates.  They keep runtime-specific facts, such
as what the execution environment reveals and which administrative steps
are ignored, out of the language definition and inside the backend
proof.

Abstract compiled layers exercise this pattern.  Audited runtimes add
erased administrative steps.  Encoded runtime examples replace
source-level booleans with a lower-level storage encoding.
Message-in-flight runtimes expose pending and delivered messages to
players and then prove either inertness, explicit cheap-talk/protocol
preservation, or a protocol-specific opacity invariant.  This last
backend is already a significant step toward the EVM target: a public
mempool with front-running is exactly pending-message visibility, which
the message-in-flight machine models and proves controllable.  A full EVM
backend --- gas, storage encoding, and contract dispatch --- and a
ConCertLean adapter would discharge the same kind of refinement
obligations at their own runtime boundary, building on this foundation.

If the backend also admits additional communication, the refinement proof
must be paired with one of the cheap-talk positions from
\Cref{sec:ref:cheap-talk}: rule the communication out of the modeled
observation space, prove it inert, or analyze a cheap-talk-extended
frontier game.  Otherwise the backend may preserve the protocol trace
law for the intended frontier strategies while still offering players
additional strategic deviations through the communication channel.

\subsection{Commit-reveal voting and governance}

Voting protocols with sealed ballots use the same commit/reveal pattern
with optional abort.  The visibility discipline of the surface language
makes the standard strategy-proofness questions of social-choice
theory~\citep{Gibbard1973,Satterthwaite1975} expressible directly: a
voting rule is \emph{strategy-proof} if truthful reporting weakly
dominates every alternative for every voter, which is the
dominance predicate on the kernel game generated from the protocol.  The
\emph{nullable yield} surface form (\Cref{sec:source:surface}) is
the protocol-level mechanism for ``a voter who refuses to reveal
contributes a $\mathsf{none}$ ballot,'' with the consequence that
later payoff expressions must explicitly handle the abstaining
voter.  Whether the protocol punishes, ignores, or rewards
abstention is a payoff-expression decision exposed by the type
system, not an implementation detail; the same is true for
threshold-of-revealers structures.  In the surface language this is a
fixed-preference protocol-game statement; full social-choice
strategy-proofness over private type spaces requires the type-space
extension discussed in \Cref{sec:disc:ext}.

\subsection{Escrow with disputes}

Two-party escrow is a small but instructive example with a
non-trivial guard.  Buyer and seller commit to delivery and
acceptance evidence; on reveal, a payoff schedule depends on
whether the evidence agrees, with a dispute branch handled by a
public coin or third-party reveal.  The non-trivial guard appears
on the dispute resolver's commit (it must select a value consistent
with prior reveals); the commit choice footprint contains only the
resolver's view.

\subsection{Coordination games}

Battle-of-the-sexes-style coordination games have multiple pure Nash
equilibria, with a mixed equilibrium quantifying the coordination
failure.  They fit directly: two simultaneous commits, two reveals,
payoff expressions encoding the coordination preference.
The Matching-Pennies example from \Cref{sec:intro:matching-pennies} uses
the same simultaneous-commit graph shape with zero-sum payoffs; changing
the payoff expression turns the same graph into a coordination game.
The kernel game exposes both the pure equilibria and, via mixed
extension, the mixed one.  Bayesian games with genuinely private
signals~\citep{Harsanyi1967} require a hidden-randomness primitive
that the surface language does not expose ($\Sample$
draws are public); this is a natural extension discussed in
\Cref{sec:disc:ext}.

\subsection{Equilibrium analysis pipelines}

The kernel-game API exposes the mathematical ingredients an
external solver adapter needs: a strategy carrier, a profile-to-
outcome distribution kernel, and a utility.  The implemented export
API packages finite \emph{pure} kernel games as player and
strategy enumerations plus kernel and utility oracles.  It is not a
standalone Gambit or PRISM-games file writer, and it does not
enumerate PMF-behavioral strategy spaces.  Those formats and libraries
are natural downstream targets rather than the artifact provided here
\citep{McKelveyMcLennanTurocy2006Gambit,
KwiatkowskaNormanParkerSantos2020PRISMGames,
LanctotEtAl2019OpenSpiel}.  The broader solver landscape includes
sequence-form and complexity results, influence-diagram models, and
regret-minimization pipelines for large imperfect-information games
\citep{KollerMegiddo1992,koller2003multi,Zinkevich2007CFR,
Bowling2015Heads-up,Brown2018Libratus}.  The canonical EFG presentation
gives a separate tree-shaped view whose payoff-vector outcome adequacy
theorem (\Cref{thm:efg}) can support future EFG exports.  The Kuhn
package
(\Cref{thm:kuhn-beh-to-mixed,thm:kuhn-mixed-to-beh,thm:kuhn-strategic-equivalence})
relates the finite mixed-pure and behavioral frontier views used before
export.  Existence corollaries for
mixed Nash, correlated equilibrium, and coarse correlated equilibrium
follow from finite-game existence results: Nash via fixed-point
theorems, and correlated or coarse correlated equilibrium by linear
programming or direct finite constructions
\citep{Nash1951,Aumann1974}.

\subsection{What this framework does not address}

The intended scope is finite, terminating, public-payoff protocols
with discrete value domains and finite-support stochastic samples.
Applications outside this scope require extensions discussed in
\Cref{sec:discussion}, including:
\begin{itemize}\itemsep1pt
\item continuous bid spaces;
\item infinite horizons or measure-theoretic kernels;
\item mechanism design with arbitrary type spaces;
\item solution concepts that depend on continuous best-response
  correspondences.
\end{itemize}
